% titel: Blatt 2, A9: Reaktionsgleichungen ausgleichen
% Teil II: die durchgerechneten Gleichungen.

\subsection*{Gegeben und gesucht}
Fünf Reaktionsgleichungen auf kleinste ganzzahlige Koeffizienten ausgleichen.

\subsection*{Lösung}
\begin{itemize}
  \item[a)] Sb$_2$S$_3$ + HCl $\to$ SbCl$_3$ + H$_2$S
        \begin{equation}
          \mathrm{Sb_2S_3} + 6\,\mathrm{HCl} \;\to\;
          2\,\mathrm{SbCl_3} + 3\,\mathrm{H_2S}.
        \end{equation}
        Bilanz: Sb $2=2$, S $3=3$, H $6=6$, Cl $6=6$.
  \item[b)] As$_2$O$_3$ + C $\to$ As + CO$_2$
        \begin{equation}
          2\,\mathrm{As_2O_3} + 3\,\mathrm{C} \;\to\;
          4\,\mathrm{As} + 3\,\mathrm{CO_2}.
        \end{equation}
        Bilanz: As $4=4$, O $6=6$, C $3=3$.
  \item[c)] Zn + H$_2$SO$_3$ + HCl $\to$ ZnCl$_2$ + H$_2$S + H$_2$O
        \begin{equation}
          3\,\mathrm{Zn} + \mathrm{H_2SO_3} + 6\,\mathrm{HCl}
          \;\to\; 3\,\mathrm{ZnCl_2} + \mathrm{H_2S} + 3\,\mathrm{H_2O}.
        \end{equation}
        Bilanz: Zn $3=3$, H $2+6=8$ vs.\ $2+6=8$, S $1=1$, O $3=3$,
        Cl $6=6$.
  \item[d)] Na$_2$HPO$_4$ + NaOH + LiCl $\to$ Li$_3$PO$_4$ + NaCl + H$_2$O
        \begin{equation}
          \mathrm{Na_2HPO_4} + \mathrm{NaOH} + 3\,\mathrm{LiCl}
          \;\to\; \mathrm{Li_3PO_4} + 3\,\mathrm{NaCl} + \mathrm{H_2O}.
        \end{equation}
        Bilanz: Na $3=3$, H $1+1=2$ vs.\ $2$, P $1=1$, O $4+1=5$ vs.\
        $4+1=5$, Li $3=3$, Cl $3=3$.
  \item[e)] C$_2$H$_2$ + O$_2$ $\to$ CO$_2$ + H$_2$O
        \begin{equation}
          2\,\mathrm{C_2H_2} + 5\,\mathrm{O_2} \;\to\;
          4\,\mathrm{CO_2} + 2\,\mathrm{H_2O}.
        \end{equation}
        Bilanz: C $4=4$, H $4=4$, O $10 = 8+2$.
\end{itemize}

\begin{ergebnis}
\begin{align*}
  \mathrm{a)\ Sb_2S_3 + 6\,HCl} &\to \mathrm{2\,SbCl_3 + 3\,H_2S},\\
  \mathrm{b)\ 2\,As_2O_3 + 3\,C} &\to \mathrm{4\,As + 3\,CO_2},\\
  \mathrm{c)\ 3\,Zn + H_2SO_3 + 6\,HCl} &\to \mathrm{3\,ZnCl_2 + H_2S + 3\,H_2O},\\
  \mathrm{d)\ Na_2HPO_4 + NaOH + 3\,LiCl} &\to \mathrm{Li_3PO_4 + 3\,NaCl + H_2O},\\
  \mathrm{e)\ 2\,C_2H_2 + 5\,O_2} &\to \mathrm{4\,CO_2 + 2\,H_2O}.
\end{align*}
\end{ergebnis}

\begin{kontrolle}
Jede Gleichung von rechts gegenprüfen: Zähle je Atomsorte beide Seiten
ab; oben ist jede Anzahl identisch (Atomzahlen in den Klammern). Alle
Koeffizienten sind die kleinsten ganzen Zahlen (kein gemeinsamer
Faktor, ggT $=1$).
\end{kontrolle}

\begin{merksatz}
Ausgleichen = Atomsorten einzeln bilanzieren und Koeffizienten als kleinste
ganze Zahlen wählen. Reine Elemente ($\mathrm{O_2}$, $\mathrm{C}$,
$\mathrm{Zn}$) zuletzt anpassen; mehrmals bilanzierte Elemente als Summe
betrachten.
\end{merksatz}
